\scanpage{154}

L'existence d'un hom can (en fait un \uline{isom.})
\[
\text{(1)}\qquad \unsure{R|K}_X^\bullet\overset{\mathbb{L}}{\boxtimes}\unsure{R|K}_Y^\bullet\xrightarrow{\sim}\unsure{R|K}_{X\times Y}^\bullet
\]
et \unsure{du fait} que l'accouplement de Verdier donne un isomorphisme
\[
\text{(2)}\qquad\deleted{\unsure{\operatorname{pr}}}F_\bullet\overset{\mathbb{L}}{\boxtimes}\unsure{\underline{D}|D}_Y G_\bullet\xrightarrow{\sim}\unsure{\underline{D}|D}_Z(\deleted{\unsure{F}}\unsure{\underline{D}|D}_X F_\bullet\overset{\mathbb{L}}{\boxtimes}G_\bullet)
\]
résulte de la formule générale (appliquée deux fois)
\margin{``\uline{Künneth local}''}
\[
\boxed{\mathbb{R}\mathcal{H}om(F_\bullet\overset{\mathbb{L}}{\boxtimes}G_\bullet,P_\bullet\overset{\mathbb{L}}{\boxtimes}Q_\bullet)\xleftarrow{\sim}\mathbb{R}\mathcal{H}om(F_\bullet,P_\bullet)\overset{\mathbb{L}}{\boxtimes}\mathbb{R}\mathcal{H}om(G_\bullet,Q_\bullet)}
\]
\inserted{$\mathbb{R}^!(f\times g)^{\unsure{!}}(P_\bullet\overset{\mathbb{L}}{\boxtimes}Q_\bullet)\simeq\mathbb{R}^!f(P_\bullet)\overset{\mathbb{L}}{\boxtimes}\mathbb{R}^!g(Q_\bullet)$}
valable si $F_\bullet,G_\bullet$ sont de tor-dim finie à ch.\ constructible, et $P_\bullet,Q_\bullet$ \deleted{à tor dim finie} \inserted{$\in\operatorname{Ob}D^{\unsure{-|+}}\deleted{\unsure{c}}$}
($A=\mathbb{Z}/n\mathbb{Z}$, $n$ premier à car $k$),
formule valable si sur les schémas \unsure{ambiants} $U,V$, \inserted{et ceux \unsure{qui sont propres} sur eux,} la résolution des singularités à la Hironaka est possible (O.K.\ en car.\ 0, et en car.\ $p$ quelcon.\ si $U,V$ de dim $\leq2$).
NB.\ Dans le cas (1), on se ramène à prendre pour $U,V$ des schémas lisses qui contiennent $X,Y$ --- \unsure{ce qui est évident dès}
\margin{153}
